% Standalone problem document, generated from the adjacent README.md.
% Compile with XeLaTeX. Mathematical content is maintained in Markdown.
\documentclass[11pt,a4paper]{article}
\usepackage[fontsize=10.5pt]{scrextend}
\usepackage[margin=22mm,headheight=15pt,headsep=9mm,footskip=11mm]{geometry}
\usepackage{fontspec}
\setmainfont{texgyrepagella}[Extension=.otf,UprightFont=*-regular,BoldFont=*-bold,ItalicFont=*-italic,BoldItalicFont=*-bolditalic]
\setsansfont{texgyreheros}[Extension=.otf,UprightFont=*-regular,BoldFont=*-bold,ItalicFont=*-italic,BoldItalicFont=*-bolditalic]
\setmonofont{texgyrecursor}[Extension=.otf,UprightFont=*-regular,BoldFont=*-bold,ItalicFont=*-italic,BoldItalicFont=*-bolditalic,Scale=MatchLowercase]
\usepackage{amsmath,amssymb}
\usepackage{unicode-math}
\setmathfont{texgyrepagella-math.otf}
\usepackage{xcolor,microtype,enumitem,needspace,fancyhdr}
\usepackage{bookmark}
\definecolor{ink}{HTML}{17344B}
\definecolor{accent}{HTML}{197D80}
\definecolor{muted}{HTML}{596773}
\hypersetup{colorlinks=true,linkcolor=accent,urlcolor=accent,pdftitle={FR-06 - Nonexistence
of a complete set of mutually unbiased bases in dimension
six},pdfauthor={Open Problems in Numerical Linear Algebra}}
\urlstyle{same}
\setlength{\parindent}{0pt}
\setlength{\parskip}{5pt plus 1pt minus 1pt}
\linespread{1.04}
\setlength{\emergencystretch}{3em}
\widowpenalty=10000
\clubpenalty=10000
\displaywidowpenalty=10000
\allowdisplaybreaks[1]
\setlist{nosep,leftmargin=1.5em}
\providecommand{\tightlist}{\setlength{\itemsep}{0pt}\setlength{\parskip}{0pt}}
\providecommand{\pandocbounded}[1]{#1}
\pagestyle{fancy}
\fancyhf{}
\fancyhead[L]{\sffamily\footnotesize\color{muted}OPEN PROBLEMS IN NUMERICAL LINEAR ALGEBRA}
\fancyhead[R]{\sffamily\bfseries\footnotesize\color{accent}FR-06}
\fancyfoot[L]{\parbox[b]{0.9\textwidth}{\sffamily\fontsize{8}{10}\selectfont\color{muted}Editorial ratings; a bounded search cannot certify that a problem remains unsolved.\\Literature check: 2026-09-08}}
\fancyfoot[R]{\sffamily\footnotesize\color{muted}\thepage}
\renewcommand{\headrulewidth}{0.3pt}
\renewcommand{\headrule}{\hbox to\headwidth{\color{accent}\leaders\hrule height \headrulewidth\hfill}}
\makeatletter
\renewcommand{\section}{\Needspace{5\baselineskip}\@startsection{section}{1}{0pt}{12pt plus 2pt minus 2pt}{4pt}{\sffamily\bfseries\large\color{ink}}}
\renewcommand{\subsection}{\Needspace{4\baselineskip}\@startsection{subsection}{2}{0pt}{9pt plus 1pt minus 1pt}{3pt}{\sffamily\bfseries\normalsize\color{ink}}}
\makeatother
\setcounter{secnumdepth}{0}
\begin{document}
{\sffamily\small\color{muted}Frames and matrix designs\par}
\vspace{3pt}
{\raggedright\hyphenpenalty=10000\exhyphenpenalty=10000\sffamily\bfseries\fontsize{21}{24}\selectfont\color{ink}Nonexistence
of a complete set of mutually unbiased bases in dimension six\par}
\vspace{7pt}
{\sffamily\small\textbf{Difficulty:} extreme\quad\textbf{Importance:} broadly
interesting\par}
\vspace{4pt}
{\color{accent}\hrule height 0.8pt}
\vspace{6pt}
\textbf{Status:} source-stated conjecture; no resolution found in
screening on 2026-09-08.

Do there fail to exist seven matrices
\(U_1,\ldots,U_7\in\mathbb C^{6\times6}\) such that \[
U_j^*U_j=I_6\quad(1\leq j\leq7),\qquad
|(U_j^*U_\ell)_{ab}|^2=\frac16
\quad(j\ne\ell,\ 1\leq a,b\leq6)?
\] Each \(U_j\) is the matrix of an orthonormal basis; the second
condition says the bases are pairwise mutually unbiased. The conjecture
is \(\operatorname{MUB}(6)<7\). It is weaker than the often discussed
assertion \(\operatorname{MUB}(6)=3\), and no claim about that sharper
value is included here.

After fixing one basis as the standard basis, the problem becomes
simultaneous feasibility of orthogonality and constant-modulus equations
for six complex Hadamard matrices. It connects structured matrix
construction, Gram matrix feasibility, and certified numerical search.
Mutually unbiased measurements also provide well-conditioned quantum
state reconstruction designs.

\subsection{References}\label{references}

\begin{enumerate}
\def\labelenumi{\arabic{enumi}.}
\tightlist
\item
  A. S. Bandeira et al., \emph{Randomstrasse 101: Open Problems of
  2025}, arXiv:2603.29571 (2026), Conjecture 22.
  \href{https://arxiv.org/html/2603.29571v1}{Paper}.
\item
  D. McNulty and S. Weigert, \emph{Mutually Unbiased Bases in Composite
  Dimensions --- A Review}, Quantum 10 (2026), article 2051;
  arXiv:2410.23997v2, revised March 26, 2026. Sections 1 and 7 state and
  survey the dimension-six existence problem.
  \href{https://arxiv.org/abs/2410.23997}{Paper}.
\item
  M. Cárdenes Wuttig and J. Tindall, \emph{A Complete Classification of
  Complex Hadamard Matrices of Order Six}, arXiv:2608.18053 (2026).
  Sections I and VI explicitly distinguish the claimed Hadamard
  classification from the still unsettled MUB problem.
  \href{https://arxiv.org/html/2608.18053}{Paper}.
\end{enumerate}

\subsection{Status check --- 2026-09-08}\label{status-check-2026-09-08}

Searched ``mutually unbiased bases dimension six solved 2026'' and ``MUB
6 Hadamard classification''; inspected the current 2026 review and
August 2026 classification. The latter explicitly states in its
conclusion that its result does not settle mutually unbiased bases in
\(\mathbb C^6\). Numerical nonexistence searches for particular families
do not exclude all seven-tuples above.
\end{document}
